\chapter{Coleta: registrar e capturar}

\eyebrow{aquisição com pareamento verificável}

\vspace{4pt}
Durante a coleta, uma carcaça é registrada antes de as imagens serem capturadas.
A imagem herda o pareamento da carcaça, e não de um nome de arquivo.

\section{Criar um lote}

Um \textbf{lote} é uma sessão de coleta (um dia de abate, um local). Ele é criado
a partir do painel do Projeto ou por importação. Cada lote agrupa as carcaças
coletadas em conjunto.

\section{Registrar uma carcaça}

O lote é selecionado e a visão \textbf{Overview} é aberta. Uma carcaça é registrada
com os seguintes campos:

\begin{center}
\begin{tabularx}{\linewidth}{@{}l X@{}}
\toprule
\textbf{Campo} & \textbf{Significado} \\
\midrule
Etiqueta física \textbf{(obrigatória)} & O identificador legível na etiqueta física.
      Ela é única por lote e restringe o erro de pareamento. \\
Animal (lab) & O identificador do animal na planilha de laboratório: o vínculo explícito. \\
Tratamento & Grupo de dieta / tratamento. \\
Estrato & Estrato de amostragem, para cobertura estratificada. \\
\bottomrule
\end{tabularx}
\end{center}

\begin{caution}[title={A etiqueta física é obrigatória}]
Uma carcaça não pode ser criada sem uma etiqueta física, e a mesma etiqueta não pode
ser usada duas vezes em um mesmo lote. Essa restrição mantém o vínculo imagem--animal
verificável e evita a falha de rótulo duplicado da coleta anterior.
\end{caution}

\section{Capturar imagens}

A carcaça é selecionada e a visão \textbf{Images} é aberta
(Figura~\ref{fig:carcass-images}). A câmera é iniciada, o ângulo de vista é escolhido
(posterior, lateral, dorsal), um quadro é capturado, a prévia é revisada e a imagem
é salva. A imagem salva é gravada na pasta da carcaça com seu \tele{carcass\_id} já
definido.

\shot{carcass-images.jpeg}{A visão Images de uma carcaça. À esquerda: captura por webcam com
um seletor de vista e uma prévia de captura; o botão \tele{Save (pair to \#IMG\_6250)}
torna o pareamento explícito. À direita, o Inspector mostra o registro completo da
carcaça selecionada: identificação, referência física, análise por modelo e
conformação.}{carcass-images}

O Inspector à direita (Figura~\ref{fig:carcass-images}) é o registro completo
da carcaça. Assim que a imagem tiver sido analisada, ele também mostra os resultados
do modelo: \tele{Fat (model)}, se o fundo foi removido, o grau
experimental e as convexidades de conformação.

\begin{note}[title={Kinect (RGB + profundidade), opcional}]
Se um Kinect e suas bibliotecas estiverem presentes, o painel de captura oferece uma
captura \tele{RGB+depth} que também registra um mapa de profundidade. Sem ele, o painel
recorre à captura por webcam e mostra \tele{Kinect unavailable (webcam only)}.
\end{note}

\section{Editar os dados da carcaça depois}

A visão \textbf{Overview} da carcaça é um formulário editável para todos os campos,
inclusive a referência física medida no animal (espessura de gordura, GR, área de olho
de lombo). Esses campos são preenchidos assim que as medidas de laboratório ficam
disponíveis, pois são os dados requeridos pela análise de correlação física
(Estação~7).
